% =============================================================
%  Pyramid 对比图 — 经典 Laplacian Pyramid vs CC-iGPT 双尺度
%  侧视立体版（强 oblique）+ bit-exact 数据流
%  编译: xelatex pyramid_compare.tex
%  嵌入 slide: \includegraphics[height=0.78\textheight]{pyramid_compare.pdf}
% =============================================================
\documentclass[tikz, border=10pt]{standalone}
\usepackage[UTF8]{ctex}
\usepackage{amsmath}
\usetikzlibrary{positioning, arrows.meta, calc,
                decorations.pathmorphing, decorations.markings}

\definecolor{ours}{HTML}{6C8CFF}
\definecolor{accent}{HTML}{FF9800}
\definecolor{rcol}{HTML}{E25555}
\definecolor{gpt}{HTML}{74C69D}

% ── oblique 偏移（侧视幅度）──
\def\ox{0.42}    % 顶面/右面水平偏移
\def\oy{0.26}    % 顶面/右面垂直偏移

% ── 3D 块宏：oblique 投影 ──
% #1 = x_center, #2 = y_bottom, #3 = width, #4 = height,
% #5 = base color, #6 = front label
\newcommand{\blockD}[6]{%
  % 地面投影
  \fill[black, opacity=0.10]
    (#1-#3/2+0.10, #2-0.12) -- (#1+#3/2+\ox+0.10, #2-0.12) --
    (#1+#3/2+\ox+0.10, #2+0.04) -- (#1-#3/2+0.10, #2+0.04) -- cycle;
  % 顶面
  \fill[#5!25, draw=#5!75!black, line width=0.4pt]
    (#1-#3/2, #2+#4) --
    (#1+#3/2, #2+#4) --
    (#1+#3/2+\ox, #2+#4+\oy) --
    (#1-#3/2+\ox, #2+#4+\oy) -- cycle;
  % 右面
  \fill[#5!70!black, draw=#5!75!black, line width=0.4pt]
    (#1+#3/2, #2) --
    (#1+#3/2+\ox, #2+\oy) --
    (#1+#3/2+\ox, #2+#4+\oy) --
    (#1+#3/2, #2+#4) -- cycle;
  % 正面
  \fill[#5!55, draw=#5!75!black, line width=0.4pt]
    (#1-#3/2, #2) rectangle (#1+#3/2, #2+#4);
  % 标签
  \node[font=\tiny, white] at (#1, #2+#4/2) {#6};
}

\begin{document}
\begin{tikzpicture}[font=\footnotesize]

  % ========== 标题 ==========
  \node[font=\bfseries\small] at (-3.0, 4.4)
    {经典 Laplacian Pyramid};
  \node[font=\scriptsize\itshape, gray!55!black] at (-3.0, 3.95)
    {Burt \& Adelson, 1983};

  \node[font=\bfseries\small] at (3.5, 4.4)
    {CC-iGPT 双尺度（本工作）};
  \node[font=\scriptsize\itshape, gray!55!black] at (3.5, 3.95)
    {coarse-conditioned AR + bit-exact ctx};

  % ========== 中间分隔虚线 ==========
  \draw[gray!35, dashed, line width=0.4pt] (0, -2.2) -- (0, 4.7);

  % ========== 左半：5 层 3D 金字塔（侧视）==========
  % 从底到顶；中间层不再标数字，只保留视觉层次
  \blockD{-3.0}{0.00}{3.20}{0.50}{ours}{$32{\times}32$ RGB}
  \blockD{-3.0}{0.68}{2.40}{0.42}{ours}{}
  \blockD{-3.0}{1.28}{1.70}{0.36}{ours}{}
  \blockD{-3.0}{1.82}{1.10}{0.32}{ours}{}
  \blockD{-3.0}{2.32}{0.55}{0.30}{ours}{$2{\times}2$}

  % 左侧侧标：用一道大括号 + 说明
  \draw[gray!55, line width=0.4pt]
    (-1.05, 0.00) -- (-0.85, 0.00) -- (-0.85, 2.62) -- (-1.05, 2.62);
  \node[font=\scriptsize, gray!50!black, anchor=west, align=left]
    at (-0.85, 1.31) {\itshape 4--5 层渐进下采样};

  % 左侧特点
  \node[font=\scriptsize, gray!50!black, align=left, anchor=north west]
    at (-5.7, -0.4)
    {$\bullet$ 多层渐进下采样（典型 4--5 层）\\
     $\bullet$ 每层独立的完整 RGB 表示\\
     $\bullet$ 通用图像分析 / 合成框架\\
     $\bullet$ 层间为差分（Laplacian 残差）};

  % ========== 右半：2 层 3D 块 ==========
  % Fine block（底层，大块）
  \blockD{3.5}{0.00}{3.00}{0.75}{ours}{Fine: $32{\times}32$ RGB}
  \node[font=\scriptsize, anchor=west, gray!55!black] at (5.45, 0.55)
    {3072 token};
  \node[font=\scriptsize, anchor=west, gray!55!black] at (5.45, 0.22)
    {\textcolor{accent}{\textbf{95.4\%}} bpd};

  % Coarse R-only block（顶层，红色突出）
  \blockD{3.50}{2.45}{0.90}{0.58}{rcol}{\textbf{R}}
  \node[font=\scriptsize, anchor=south, rcol!50!black] at (3.5, 3.20)
    {\textit{R-only}};
  \node[font=\scriptsize, anchor=east, gray!55!black] at (3.0, 2.74)
    {$8{\times}8{\times}1$};
  \node[font=\scriptsize, anchor=west, gray!55!black] at (4.85, 2.92)
    {64 token};
  \node[font=\scriptsize, anchor=west, gray!55!black] at (4.85, 2.62)
    {\textcolor{accent}{\textbf{4.6\%}} bpd};

  % ========== bit-exact 数据流 ==========
  % 主流（snake 装饰，从 coarse 底面流到 fine 顶面）
  \draw[accent!85!black, line width=1.4pt,
        decorate, decoration={snake, amplitude=0.8pt,
                              segment length=4pt, pre length=2pt,
                              post length=4pt},
        -{Stealth[length=5pt, width=4pt]}]
    (3.50, 2.45) -- (3.50, 0.80);

  % α·ctx 块（3D 面，代替原圆角卡片）
  \blockD{3.50}{1.42}{1.30}{0.42}{gpt}{}
  % 块内文字（白色，比 \tiny 大）
  \node[font=\scriptsize\bfseries, white]
    at (3.50, 1.63) {$\alpha \cdot \mathrm{ctx}$};
  % iGPT 标签（块右侧，绿色深一档）
  \node[font=\scriptsize\bfseries, gpt!45!black, anchor=west]
    at (4.85, 1.78) {iGPT};

  % 流上沿途的 pipeline stage 小节点（右侧标注）
  \fill[accent!85!black] (3.50, 2.20) circle (1.2pt);
  \node[font=\tiny, gray!55!black, anchor=west] at (3.62, 2.20)
    {\texttt{/255 dequant}};

  \fill[accent!85!black] (3.50, 2.00) circle (1.2pt);
  \node[font=\tiny, gray!55!black, anchor=west] at (3.62, 2.00)
    {\texttt{bilinear UP $32{\times}32$}};

  \fill[accent!85!black] (3.50, 1.20) circle (1.2pt);
  \node[font=\tiny, gray!55!black, anchor=west] at (3.62, 1.20)
    {\texttt{re-tokenize}};

  \fill[accent!85!black] (3.50, 1.00) circle (1.2pt);
  \node[font=\tiny, gray!55!black, anchor=west] at (3.62, 1.00)
    {\texttt{fine.token\_embed}};

  % 流右侧的总标签：bit-exact stream
  \node[font=\scriptsize\itshape, accent!70!black, anchor=south,
        rotate=90]
    at (3.20, 1.62) {bit-exact stream};

  % 右侧特点
  \node[font=\scriptsize, gray!50!black, align=left, anchor=north west]
    at (1.6, -0.4)
    {$\bullet$ 仅 \textbf{2 层}（coarse + fine）\\
     $\bullet$ Coarse 仅 R 通道（$\sim$2\% 预算）\\
     $\bullet$ Fine 通过 sub-pixel AR 自学 G/B\\
     $\bullet$ Coarse$\to$fine 经 \emph{bit-exact} 注入\\
     $\bullet$ 仅引入 1 个标量参数 $\alpha$};

\end{tikzpicture}
\end{document}
